% !TeX root = ../CourseReport.tex
%=============================================================================
%  PART 3 --- USER ACCEPTANCE TESTING (UAT)
%  PAGE BUDGET: 3 pages (excluding title, ToC)          Rubric weight: 10 pts
%
%  NOTE FOR THE TEAM: the cohort composition, the seven Likert means, and the
%  free-text themes below are the figures reported in the final presentation.
%  Search for "CONFIRM" before submission.
%=============================================================================
\chapter{User Acceptance Testing}
\label{ch:uat}

\section{User Case Study: Who We Studied and Why}

Dishify's target user is the \textbf{everyday home cook who already owns food
but has no plan}: someone who cooks for themselves or a small household on most
days, shops irregularly, and repeatedly faces the same 6~p.m. question --- what
can I make with what is in the fridge right now?

We chose this profile deliberately over two alternatives. \emph{Novice cooks}
would have stressed instruction quality more than recommendation quality, and
\emph{enthusiast cooks} already own the mental index Dishify is meant to
replace, so neither would exercise the core pantry-to-recipe loop. The everyday
cook is also the group for whom our two value claims are testable: they suffer
decision fatigue, and their unused ingredients become waste.

Recruitment was therefore screened on \emph{cooking frequency} rather than
skill. Nine participants were drawn from the team's wider student and personal
network, in two deliberately different segments. \textbf{Four cook daily} ---
the group with the highest exposure to repeat decision fatigue, who would show
whether suggestions stay varied and whether a standing pantry is convenient to
maintain. \textbf{Five cook weekly} --- planning in batches from a longer-lived
pantry, who would stress ingredient coverage, tolerance for missing items, and
the usefulness of the shopping gap. Their self-reported pain points confirmed
the premise before they saw the product: \textbf{5 of 9} named ``I don't know
what to cook'' and \textbf{4 of 9} ``recipes take too long''.

The trade-off of this cohort is worth stating. It gave us participants who
genuinely have the problem and would speak candidly, but a convenience sample
is small ($n=9$), student-skewed, and prone to politeness bias --- which is why
we weight the behavioural and free-text findings at least as heavily as the
Likert means.

\section{Survey Design}

The study combined a \textbf{task-based usability test} with a
\textbf{structured survey}, so that what people did could be compared with what
they said. Participants were given the deployed web application and a short
scenario --- \emph{cook something tonight from what you actually have at
home} --- and were asked to think aloud while working through a fixed task
sequence:

\begin{enumerate}[itemsep=1pt, topsep=3pt]
    \item register an account and complete first-time setup;
    \item record any real allergies, diets, or restrictions in
          \emph{Preferences};
    \item enter a pantry reflecting their own kitchen, using typing, voice, or
          the photo scan;
    \item request recommendations and judge the ranked list;
    \item open one recipe and decide whether they would actually cook it.
\end{enumerate}

Each task carried an observed success criterion --- completed unaided, with a
hint, or abandoned --- and we noted where participants hesitated. The survey
then asked for agreement on a five-point Likert scale across seven dimensions
chosen to map onto the project's value claims rather than to generic
usability: \textbf{ease of use} and \textbf{recommendation process} (is the
pantry-first flow learnable without instruction?), \textbf{recipe relevance}
(are suggestions actually cookable from the stated pantry?),
\textbf{explanation clarity} (do ``why it fits'' and ``watch for'' carry real
information?), \textbf{allergy and diet trust} (would the user rely on the hard
filters for a restriction that matters to them?), \textbf{substitution quality}
(is the handling of missing ingredients helpful?), and \textbf{likelihood to
reuse} --- our retention proxy, and the hardest question we asked. Two
free-text questions closed the survey: what worked well, and what would have to
change before they would use Dishify again next week.

\section{Survey Execution}

All nine participants completed the full task sequence and the survey on the
deployed instance (\url{http://167.233.165.44}), each using their own real
pantry and real restrictions rather than a scripted example. Sessions ran
individually; the observer recorded task outcomes and think-aloud remarks,
after which the participant filled in the survey unobserved to limit courtesy
effects. Every participant reached a ranked result list and opened at least one
recipe, so no session was excluded, and all seven Likert dimensions received
nine responses. % CONFIRM exact session dates and observer per session.

%-----------------------------------------------------------------------------
%  TEAM ACTION (biggest single scoring gap identified in review):
%  The rubric rewards *auditable* UAT evidence. We currently report only the
%  aggregate means. If you still have the survey export, uncomment and fill the
%  table below (it fits if you shorten the "Survey Design" paragraph slightly),
%  and attach the raw CSV to the Moodle submission.
%
%  \begin{table}[!ht]
%      \centering\footnotesize
%      \caption{Per-participant task outcomes and ratings ($n=9$).}
%      \label{tab:uat-raw}
%      \begin{tabular}{@{}clccccccc@{}}
%          \toprule
%          \textbf{P} & \textbf{Cooks} & \textbf{Tasks done} & \textbf{Ease} &
%          \textbf{Proc.} & \textbf{Trust} & \textbf{Clar.} & \textbf{Rel.} &
%          \textbf{Reuse} \\
%          \midrule
%          P1 & daily  & 5/5 & 5 & 5 & 4 & 4 & 4 & 4 \\
%          ... one row per participant ...
%          \bottomrule
%      \end{tabular}
%  \end{table}
%
%  Also add: session dates, and 3--5 short anonymised verbatim quotes.
%-----------------------------------------------------------------------------

\section{Result Analysis and Evaluation}
\label{sec:uat-analysis}

\begin{figure}[!ht]
    \centering
    \begin{tikzpicture}
        \begin{axis}[
            xbar,
            width=0.74\textwidth,
            height=4.6cm,
            xmin=0, xmax=5,
            xtick={0,1,2,3,4,5},
            enlarge y limits=0.12,
            bar width=11pt,
            axis lines*=left,
            xmajorgrids, grid style={draw=black!12},
            tick align=outside,
            xlabel={mean agreement (1--5), $n=9$},
            symbolic y coords={Reuse,Substitution,Relevance,Clarity,Trust,Process,Ease},
            ytick=data,
            nodes near coords={\pgfmathprintnumber[fixed, precision=2]{\pgfplotspointmeta}},
            nodes near coords style={font=\footnotesize, anchor=west},
            every node near coord/.append style={xshift=1pt},
        ]
        \addplot[fill=dishifygreen, draw=black!40] coordinates {
            (4.67,Ease) (4.56,Process) (4.44,Trust)
        };
        \addplot[fill=dishifyorange!35, draw=black!40] coordinates {
            (4.11,Clarity) (4.00,Relevance)
        };
        \addplot[fill=dishifyorange!75, draw=black!40] coordinates {
            (3.88,Substitution) (3.67,Reuse)
        };
        \end{axis}
    \end{tikzpicture}
    \caption[User acceptance results]{Mean Likert agreement across the seven
    evaluated dimensions ($n=9$). Interaction scores (top, green) are strong;
    the content-quality and retention scores (bottom) are where the product is
    weakest.}
    \label{fig:uat-scores}
\end{figure}

Figure~\ref{fig:uat-scores} shows a clear split. Everything about
\emph{operating} Dishify scored highly --- ease of use $4.67$, recommendation
process $4.56$ --- so the pantry-first model is understood without instruction,
and our central bet on inverting the usual dish-first search was correct.
Allergy and diet trust at $4.44$ is the result we cared most about:
participants with real restrictions were willing to rely on the hard filters,
validating the decision to enforce restrictions in retrieval rather than in the
prompt.

The weaker scores concern \emph{content}, not interaction, and are consistent
with each other. Recipe relevance ($4.00$), substitution quality ($3.88$), and
likelihood to reuse ($3.67$) form a descending chain: recipes that are only
approximately cookable produce weak substitution advice, and weak substitution
advice erodes the reason to come back. Explanation clarity ($4.11$) sits
between --- the explanations are trusted, but were found terse.

The free-text answers explain the numbers. What users valued matched our value
claims almost exactly: Dishify \emph{removes the decision}, \emph{uses up food
they already own}, and the explanations \emph{make the ranking legible} rather
than arbitrary. What blocked reuse was concrete: generation felt slow;
ingredient matching was too literal (``spring onion'' not matching ``onion'');
instructions needed more detail; substitutions were weak; and onboarding did
not make the value obvious quickly enough. Participants also asked most often
for bookmarking, cooking-time and difficulty filters, health-goal filters,
recipe photos, and a print-friendly view.

\paragraph{What we incorporated, and what we did not.}
Table~\ref{tab:uat-actions} records the disposition of each finding. Two were
fixed in the time remaining; literal matching was partially addressed; latency
we measured rather than guessed, which redirected the fix from retrieval to the
explanation step; and the convenience requests are recorded as future work
rather than quietly dropped. The honest summary is that Dishify passed
acceptance on \emph{usability and trust} and did not yet pass on
\emph{retention}: a $3.67$ likelihood to reuse from a friendly cohort is a
warning, not a pass mark, and it is the number we would optimise first in any
continuation.

\begin{table}[!htbp]
    \centering
    \footnotesize
    \caption{Disposition of the main UAT findings.}
    \label{tab:uat-actions}
    \begin{tabularx}{\textwidth}{@{}p{36mm}p{22mm}X@{}}
        \toprule
        \textbf{Finding} & \textbf{Status} & \textbf{Action taken / rationale} \\
        \midrule
        Instructions too terse & Incorporated & LLM step augmentation rewrites terse directions into timed, numbered steps. \\
        Explanations too short & Incorporated & Reasoning prompt tightened to produce specific, recipe-grounded bullets. \\
        Matching too literal & Partial & Plural normalisation and token-containment matching added; synonym ontology deferred. \\
        Generation too slow & Diagnosed & Explanation is $\approx$\,47\,\% of median response time; fix is streaming/deferred explanation. \\
        Weak substitutions & Deferred & Needs an ingredient-substitution knowledge base. \\
        Onboarding, bookmarks, filters, photos, print view & Deferred & Valuable but non-core; recorded as future work (Section~\ref{sec:pm-progress}). \\
        \bottomrule
    \end{tabularx}
\end{table}
